% Incorporación al final de extension.tex; no contiene una sección principal.
% Procedencia: RESULTADO_RECUPERADO; explicitación sectorial de las hipótesis.
% Fuentes leídas: monografía Doble proyección holográfica (cantera editorial),
% manuscrito/local/01_tesis_recuperada.tex y
% manuscrito/generated/algebra_holonomica.tex, especialmente §§1--7 y 21.
% Definición rectora de número: manuscrito/flattened/main_autosuficiente.tex,
% secciones que comienzan en líneas 22137 y 23000 de la cantera de 775 páginas.
% No se adopta su presentación global como prueba de todas las flechas TPK.

\subsection{El número como sección compatible y el valor como lectura}
\label{hist:seccion}

La distinción entre estado y observación permite precisar qué objeto se
prolonga cuando aumenta la resolución. A profundidad \(n\), sea \(X_n\)
el conjunto de estados admisibles producidos por APP--TRIT--TPK. Sus datos
incluyen posición, hojas suma y producto, residuos y cocientes, régimen,
orientación, acarreo, fase, ruta, frontera, memoria y supervivencia. No son
coordenadas independientes: las evaluaciones aritméticas y los transportes
anteriores imponen sus relaciones. Una prolongación conserva esas relaciones
en cada antecedente, aunque añada nuevos pasos y nuevos datos de memoria.
Éste es el sentido de la aritmética de historias desarrollada en
(secciones~\ref{sec:extension} y~\ref{mono-ampl:conexion}).

Sean \(r_{m,n}:X_m\to X_n\), para \(m\geq n\), las restricciones de
profundidad, con \(r_{n,n}=\operatorname{id}\) y
\(r_{\ell,n}=r_{m,n}\circ r_{\ell,m}\). Cada restricción conserva el
prefijo y la información ya determinada en él.

\begin{definition}[Sección numérica compatible]
Un número HMT es una historia compatible del sistema de estados:
\begin{equation}
 x=(x_n)_{n\geq0}\in X_\infty:=\varprojlim_n X_n,
 \qquad r_{m,n}(x_m)=x_n.
 \label{hist:limite-inverso}
\end{equation}
Un lector es una operación adicional sobre un dominio especificado de estas
historias. Una cifra es una observación finita; un valor escalar es una
evaluación del lector. El par formado por la sección y su lector constituye
una presentación leída, no una redefinición de la identidad de la sección.
\end{definition}

En el dominio arquimediano \(X_\infty^{\mathrm{Arch}}\), el lector asocia
a cada prefijo un cilindro cerrado, conserva su encajamiento y hace tender
su diámetro a cero. Así determina
\begin{equation}
 \operatorname{ev}_{\mathbb R}:X_\infty^{\mathrm{Arch}}\longrightarrow
 \mathbb R,\qquad
 \{\operatorname{ev}_{\mathbb R}(x)\}=\bigcap_n I_n(x).
 \label{hist:evaluacion}
\end{equation}
La construcción concreta de estos cilindros y la decisión de sus fronteras
se desarrollan en la sección~\ref{sec:generacion}. La definición
\eqref{hist:limite-inverso} no afirma por sí sola que la imagen de
\eqref{hist:evaluacion} sea toda la recta: una afirmación de cobertura
requiere su propio teorema. Tampoco una igualdad de valores establece, sin
un resultado de inyectividad, la igualdad de sus historias. Esta separación
permite hablar de precisión como profundidad conservando la diferencia
entre el antecedente aritmético y su realización escalar.

\subsection{Composición de historias y multiplicador de memoria}
\label{hist:algebra}

Fijemos un sector de estados a profundidad \(n\) y una categoría
\(\mathcal C\) cuyas flechas sean sus historias admisibles. La composición
se escribe \(vu=v\circ u\): primero se recorre \(u\) y después \(v\).
Los pasos elementales proceden de la selección, el transporte y la
actualización del TPK. Una fibra trítica local admite la escritura
\(\mathbb R[J_x]/(J_x^2+\tau(x)1_x)\), donde
\(\tau(x)\in\{+1,0,-1\}\). El régimen y la orientación pertenecen al
estado que se transporta; una transición entre regímenes no equivale a
identificar sus álgebras cuadráticas.

La siguiente formulación precisa un sector algebraico de esa composición.
Para cada objeto \(x\), se da un álgebra real asociativa unitaria \(A_x\),
que puede ser no conmutativa. Para cada \(u:x\to y\), se da un
homomorfismo unital \(\rho_u:A_x\to A_y\). Suponemos una acción estricta:
\begin{equation}
 \rho_{\operatorname{id}_x}=\operatorname{id}_{A_x},\qquad
 \rho_v\rho_u=\rho_{vu}.
 \label{hist:accion-estricta}
\end{equation}
Para cada pareja componible \(u:x\to y\), \(v:y\to z\), fijamos un
multiplicador central invertible
\(\Omega(v,u)\in Z(A_z)^\times\). Se exige la normalización
\(\Omega(\operatorname{id}_y,u)=\Omega(u,\operatorname{id}_x)=1_y\)
y, para cada triple componible, la identidad
\begin{equation}
 \rho_w\bigl(\Omega(v,u)\bigr)\Omega(w,vu)
   =\Omega(w,v)\Omega(wv,u).
 \label{hist:cociclo}
\end{equation}
Son hipótesis explícitas sobre los transportes y sus coeficientes, no una
consecuencia de llamar holonómica a la dinámica. Al aplicar esta presentación
a un sector TPK deben identificarse sus \(\rho_u\) y \(\Omega(v,u)\).
Las transiciones expresadas mediante bimódulos requieren conservar esa
estructura y no quedan reemplazadas por \eqref{hist:accion-estricta}.

El espacio de sumas finitas
\[
 \mathcal H(\mathcal C,A,\rho,\Omega)
   =\bigoplus_{u\in\operatorname{Mor}(\mathcal C)} A_{t(u)}T_u
\]
recibe el producto bilineal
\begin{equation}
 (aT_v)(bT_u)=a\rho_v(b)\Omega(v,u)T_{vu},
 \label{hist:producto}
\end{equation}
si las flechas son componibles, y cero en otro caso. Aquí \(t(u)\)
designa el estado terminal. La suma reúne historias con coeficientes; el
producto concatena historias conservando el multiplicador del registro.

\begin{proposition}[Asociatividad sectorial con multiplicador central]
Bajo \eqref{hist:accion-estricta} y \eqref{hist:cociclo}, el producto
\eqref{hist:producto} es asociativo. Si el conjunto de objetos es finito,
su unidad es \(\sum_x1_xT_{\operatorname{id}_x}\).
\end{proposition}
\begin{proof}
Para \(u:x\to y\), \(v:y\to z\), \(w:z\to t\), sean
\(a\in A_y\), \(b\in A_z\), \(c\in A_t\). Los dos coeficientes
de \(T_{wvu}\) son
\begin{align*}
 ((cT_w)(bT_v))(aT_u):\quad&
 c\rho_w(b)\Omega(w,v)\rho_{wv}(a)\Omega(wv,u),\\
 (cT_w)((bT_v)(aT_u)):\quad&
 c\rho_w(b)\rho_w\rho_v(a)\rho_w(\Omega(v,u))\Omega(w,vu).
\end{align*}
La acción estricta iguala \(\rho_w\rho_v(a)\) con \(\rho_{wv}(a)\).
La centralidad de \(\Omega(w,v)\) permite desplazarlo a la derecha de
ese factor; \eqref{hist:cociclo} iguala entonces los productos restantes.
No se permutan los coeficientes \(a,b,c\). Si el triple no es componible,
ambas parentizaciones son cero por sus estados iniciales y terminales.
La bilinealidad concluye la prueba. Las identidades locales y la
normalización de \(\Omega\) verifican la unidad indicada.
\end{proof}

El acarreo del calendario proporciona un ejemplo efectivo. En la categoría
de un objeto y flechas \(C_9\), tomemos coeficientes
\(A=\mathbb R[z,z^{-1}]\), acción identidad y
\[
 \Omega(b,a)=z^{c_0(a,b)},\qquad
 c_0(a,b)=\left\lfloor\frac{a+b}{9}\right\rfloor,
 \quad 0\leq a,b<9.
\]
Las dos sumas de acarreos de un triple son
\(\lfloor(a+b+c)/9\rfloor\); por ello se cumple
\eqref{hist:cociclo}. La variable formal \(z\) registra la vuelta sin
asignarle un valor numérico. Imponer después \(z^{12}=1\) recupera el
registro dodecafásico del calendario \eqref{eq:extension108}; conservar
\(z\) sin esa relación distingue vueltas sucesivas. El ejemplo realiza
la memoria de este calendario, sin identificarla con la totalidad de la
memoria TPK.

\subsection{Doble varianza del refinamiento y conservación de la unidad}
\label{hist:doble-varianza}

Al restringir estados se prolongan observables en sentido contrario.
Consideremos las álgebras de funciones, con producto puntual,
\(D_n=\operatorname{Fun}(X_n,\mathbb R)\). La precomposición define
\begin{equation}
 r_{m,n}^*:D_n\longrightarrow D_m,
 \qquad r_{m,n}^*f=f\circ r_{m,n}.
 \label{hist:precomposicion}
\end{equation}
Son morfismos unitales; son además inyectivos cuando las restricciones
correspondientes son sobreyectivas. El límite directo
\(D_{\mathrm{cyl}}=\varinjlim_n(D_n,r_{n+1,n}^*)\) reúne las
observaciones de profundidad finita. No se identifica aquí este álgebra
con toda el álgebra de rutas: prolongar los operadores de transporte exige
también sus compatibilidades con el refinamiento.

\begin{lemma}[Unidad y naturalidad de la evaluación]
Si \(e_x\) es el indicador de \(x\in X_n\), se cumple
\begin{equation}
 r_{m,n}^*e_x=\sum_{y:\,r_{m,n}(y)=x}e_y,
 \qquad r_{m,n}^*1_n=1_m,
 \label{hist:unidad}
\end{equation}
y, para todo \(f\in D_n\), \(y\in X_m\),
\[
 \langle r_{m,n}^*f,y\rangle=\langle f,r_{m,n}y\rangle.
\]
Cada sección \(x\in X_\infty\) determina una evaluación bien definida
\(D_{\mathrm{cyl}}\to\mathbb R\), dada por \([f]\mapsto f(x_n)\)
para \(f\in D_n\).
\end{lemma}
\begin{proof}
Las fibras de \(r_{m,n}\) forman una partición de \(X_m\). Sus
indicadores son ortogonales y suman \(1_m\), lo que prueba
\eqref{hist:unidad}. La naturalidad es la definición de precomposición.
Finalmente, \((r_{m,n}^*f)(x_m)=f(x_n)\), de modo que dos representantes
compatibles dan la misma evaluación en el límite directo.
\end{proof}

La unidad global es la función constante uno sobre todas las posibilidades;
la sección individual selecciona una historia compatible. Conservar la
primera no significa inmovilizar la segunda. El refinamiento puede aumentar
la cantidad de extensiones y la profundidad de cada historia manteniendo
la unidad y todas las evaluaciones de sus antecedentes.

\subsection{Retorno de fase, representación de Weyl y dos lecturas}
\label{hist:puente-lecturas}

La ley \eqref{eq:holonomia-memoria} adquiere una realización operatoria
en la suspensión bilateral que se construirá en la
sección~\ref{sec:generacion}. Allí \(T\) es el avance de un paso,
\(V_9\) observa la fase nonádica, \(W_\delta\) la fase de rotación y
\(D_M\) el grado de memoria. Escribiendo
\(\widehat\Gamma_9=T^9\) para la vuelta en esta representación,
las relaciones de Weyl \eqref{eq:weyl-relojes} implican, sobre vectores
de soporte finito,
\begin{equation}
 \begin{aligned}
 V_9\widehat\Gamma_9&=\widehat\Gamma_9V_9,\qquad
 W_\delta\widehat\Gamma_9=e^{2\pi i9\delta}
     \widehat\Gamma_9W_\delta,\\
 [D_M,\widehat\Gamma_9]&=\widehat\Gamma_9.
 \end{aligned}
 \label{hist:tres-retornos}
\end{equation}
La primera igualdad sigue de \(\zeta_9^9=1\); la segunda, de nueve
aplicaciones de la covariancia de \(W_\delta\); la tercera expresa
el incremento unitario de memoria. La fase finita retorna, mientras la
rotación irracional y la memoria distinguen el estado siguiente. Los
caracteres y el parámetro \(\delta\) se realizan después desde la
construcción numérica allí especificada: no seleccionan retrospectivamente
las semillas. Esta representación describe observables y traslaciones de
la suspensión; no reemplaza todas las coordenadas ni todos los operadores
del TPK.

La doble lectura prolonga la misma distinción entre historia y observable.
En su dominio arquimediano se aplica \eqref{hist:evaluacion}; en el
dominio incidencial calibrado, los mapas \(Q_n\) de la
sección~\ref{exc:doble-lectura} conservan las palabras y los marcos de
frontera. La igualdad allí demostrada
\(r_{n,m}Q_n=Q_mp_{n,m}\) determina su prolongación \(Q_\infty\).
Sobre el dominio común de ambas lecturas pueden considerarse conjuntamente
\(\operatorname{ev}_{\mathbb R}(x)\) y \(Q_\infty(x)\): la
primera retiene el valor; la segunda, la incidencia transportada que su
codominio conserva. La sección antecede a ambas. Esta articulación permite
pasar de la aritmética de historias a las constantes y a la lectura
excepcional sin convertir sus realizaciones en nuevos datos iniciales.
